% 解法の選択を、計算の前に確認する短い練習。
\subsection{操作を選ぶ練習}\label{節：操作を選ぶ練習}
  ここでは\bookterm{general}{一般項}まで一気に計算する前に,新しい量を選び,その更新だけを確かめる.
  引く・割る・\bookterm{logarithm}{対数}を取るという操作が,それぞれ何を簡単にするかを比較しよう.

  \noindent\bookblocklink{operation-choice}{problem}{answer}{【問題】}\par\nobreak
  \begin{enumerate}[format=\bookitemlink{operation-choice}{problem}{answer}]
    \item $a_1=1$,\,$a_{n+1}=3a_n+2n+1$とする.
          定数を引くだけで付加項を消せるか.
          $g(n)=un+v$を引く場合の$u,v$を求め,\,$b_n=a_n-g(n)$の更新式と初項を示せ.
    \item $a_1=2$,\,$a_{n+1}=\frac{n+2}{n}a_n+(n+1)(n+2)$とする.
          $b_n=\frac{a_n}{n}$と$c_n=\frac{a_n}{n(n+1)}$を候補とし,それぞれの更新式を求めよ.
          どちらで$a_n$の係数を$1$にそろえられるか.
    \item $a_1=-2$,\,$a_{n+1}=2^{2n+1}a_n$とする.
          $b_n=\log_2a_n$は使えるか.
          \bookterm{logarithm}{対数}を使わず,\,$c_n=\frac{a_n}{2^{n^2}}$の更新式と初項を求めよ.
  \end{enumerate}

  第2問は,\sectionref*{節：階比型}で倍率をそろえた操作を,付加項がある場合に使う練習である.
  一般に
  \[
    a_{n+1}=p_na_n+q_n
  \]
  に対して,既知の非零の\bookterm{sequence}{数列}$h_n$が$h_{n+1}=p_nh_n$を満たすなら,
  \[
    b_n=\frac{a_n}{h_n},\qquad
    b_{n+1}-b_n=\frac{q_n}{h_{n+1}}
  \]
  となる.倍率をそろえる操作の後に,\bookterm{difference}{階差}を足し戻す操作が続く.
  $p_n\ne0$なら,\,$h_1=1$,\,$h_n=\prod_{k=1}^{n-1}p_k$を候補にできる（空積は$1$）.
  係数が零になる場合はこの割り算を使わず,その段階の更新を別に確認する.
  また,残った和がいつでも簡単に計算できるとは限らない.

  \noindent\bookblocklink{operation-choice}{answer}{problem}{【解答と選択の理由】}\par\nobreak
  \begin{enumerate}[format=\bookitemlink{operation-choice}{answer}{problem}]
    \item $a_n-\alpha$では付加項の$2n$が残る.
          $g(n+1)=3g(n)+2n+1$から
          \[
            u=3u+2,\qquad u+v=3v+1
          \]
          となり,\,$u=v=-1$である.
          従って$b_n=a_n+n+1$により$b_{n+1}=3b_n$,\,$b_1=3$を得る.
          引く量を定数から一次式へ変えたことで,付加項を消せた.

    \item 候補ごとに一歩だけ更新すると,
          \[
            b_{n+1}=\frac{n+2}{n+1}b_n+n+2
          \]
          と
          \[
            c_{n+1}=c_n+1,\qquad c_1=1
          \]
          を得る.
          $h_n=n(n+1)$なら
          \[
            \frac{h_{n+1}}{h_n}=\frac{n+2}{n}
          \]
          が元の倍率と一致する.
          この一致を手掛かりに割る因子を選ぶ.
          $n\ge1$なので分母は非零であり,\,$c_n=n$から$a_n=n^2(n+1)$も得られる.

    \item 正の倍率で更新するので,各項は負であり,\,$\log_2a_n$は実数として定義できない.
          一方,\,$2n+1=(n+1)^2-n^2$から
          \[
            c_{n+1}=c_n,\qquad c_1=-1
          \]
          となる.従って$a_n=-2^{n^2}$を得る.
          \bookterm{normalize}{正規化}で割る$2^{n^2}$は常に非零であり,元の項が正である必要はない.
          初項が正の比較例は\bookproblemref{practice-basic}{42}{標準問42}へ.
  \end{enumerate}
